\documentclass[palette=none,mode=light]{impaTeXPoster}

\usepackage[utf8]{inputenc}
\usepackage[T1]{fontenc}
\usepackage{lipsum}

\title{Poster Teste}
\subtitle{Exemplo de subtítulo}
\author{Bruno Pereira, Maria Silva, João Souza}
\authoremail{bruno@impa.br, maria@impa.br, joao@impa.br}
\posterlogo{images/impatech.png}    % logo à esquerda
\rightposterlogo{images/tex.png}    % logo à direita

\begin{document}
\maketitle

\useblockstyle{testStyle}
\begin{columns}

% ==========================================
% LEFT COLUMN
% ==========================================
\column{0.5}

\begin{postercard}[yellow50, yellow50, black, yellow400,black]{Introdução}
Este é um exemplo de seção usando a classe ImpaTeXPoster.
\lipsum[1]
\end{postercard}


\begin{postercard}[pink50, pink50, black, pink400,black]{Objetivo}
Mostrar como usar diferentes cores no ambiente.

\lipsum[2]
\end{postercard}

\block{Definition}{
A space $X$ is compact if every open cover admits a finite subcover.
}

\block{}{
The notion of compactness evolved during the nineteenth century.

\lipsum[1-2]
}

% ==========================================
% RIGHT COLUMN
% ==========================================
\column{0.5}

\block{Proposition}{
Every compact subset of a Hausdorff space is closed.

\lipsum[3]
}


\block{Corollary}{
Compact subsets of \(\mathbb{R}^n\) are closed and bounded.
}

\block{}{
    Proof. \\
    \lipsum[3]

}

\begin{postercard}[green50, green50, black, green400,black]{Resultados}
\lipsum[3]
\end{postercard}

% primeiro parâmetro é a cor do Fundo do título
% segundo parâmetro é a cor do do testo
% terceiro parâmetro é cor do texto
% quarto parâmetro é a borda do postercard

\begin{postercard}[orange50, orange50, black, orange400,black]{Metodologia}
\lipsum[4]
\end{postercard}

\end{columns}

\begin{columns}

%Left Column
\column{0.33}

\begin{postercard}[yellow50, yellow50, black, yellow400,black]{Metodologia}
\lipsum[5]
\end{postercard}

%Middle Column
\column{0.33}

\begin{postercard}[pink50, pink50, black, pink400,black]{Metodologia}
\lipsum[6]
\end{postercard}


%Right Column
\column{0.33}

\begin{postercard}[green50, green50, black, green400,black]{Metodologia}
\lipsum[7]
\end{postercard}

\end{columns}


\posterfooter{https://github.com/impaTeX/modelo-cartaz/}{%

    \textbf{{\LARGE Referências:}}

    \begin{thebibliography}{99}
    \bibitem{popper}
    POPPER, Karl.
    \textit{A lógica da pesquisa científica}.
    São Paulo: Cultrix, 2002.

    \bibitem{chalmers}
    CHALMERS, A. F.
    \textit{O que é ciência, afinal?}
    São Paulo: Brasiliense, 1999.
    \end{thebibliography}
}

\end{document}